% Source: https://artofproblemsolving.com/community/c3h2779431_2016_mbmt_team_round__montgomery_blair_math_tournament
% Stable identity: aops-topic-2779431-q113
The midpoint theorem makes $MNPQ$ a parallelogram with $NP=\dfrac{AC}{2}=10$. Its diagonals bisect at $X$, so triangle $NXP$ has sides 12,10,10. Taking $NX=12$ as base, its altitude is $\sqrt{10^2-6^2}=8$, giving area 48. The four triangles around $X$ have equal area, so $[MNPQ]=192$. A midpoint parallelogram has half the original quadrilateral’s area (subtract the four corner triangles, or apply the midpoint theorem on diagonal $AC$). Therefore $\boxed{[ABCD]=384}$.
\begin{asy}
usepackage("amsmath");
size(7cm);
pair A=(6,-8),B=(18,8),C=(-6,8),D=(-18,-8),M=(-6,-8),Nn=(12,0),P=(6,8),Q=(-12,0),X=(0,0);draw(A--B--C--D--cycle);draw(M--Nn--P--Q--cycle,blue);draw(M--P,dashed);draw(Nn--Q,dashed);draw(A--C,gray);label("$A$",A,S);label("$B$",B,E);label("$C$",C,N);label("$D$",D,W);label("$M$",M,S);label("$N$",Nn,E);label("$P$",P,N);label("$Q$",Q,W);label("$X$",X,S);
\end{asy}
